\section{Marking ambiguities}
\label{app:marking-ambiguities}

The relative formulation uses the following complete action of the auxiliary
choices.

\begin{center}
\begin{tabular}{@{}l p{0.65\linewidth}@{}}
\hline
operation & effect \\
\hline
axis switching &
\(C\mapsto DCD\), while every triangle product and \(Z\) is fixed; this is
the equivalence already built into the marked source. \\
axis relabelling &
simultaneously permutes the rows and columns of \(C\) and the variables of
\(Z\); the proposition transports equivariantly and does not choose an order. \\
five-label relabelling &
simultaneously permutes the \(q_i\), the coordinates \(y_i\), and the
two-subset labels; \(\sigma_3\) is fixed and \(\beta\) is equivariant.
Relabelling only one side changes \(\mathfrak m\). \\
chart scaling &
rescales the linear identification and the transported cubic generator while
leaving the projective chart fixed.  It is not quotiented: \(q_1=xyz\) fixes
the normalization.  In particular, scaling by \(-1\) reverses the generator. \\
golden Galois conjugation &
after transport by \(R\), sends \(C\), \(Z\), and the normalized chart lift
to their negatives at the golden fibre. \\
incidence deck exchange &
interchanges \(B_+\) and \(B_-\) and negates the odd generator; relative to
\(\mathfrak m\), the attached source and harmonic generator are negated by
definition. \\
\hline
\end{tabular}
\end{center}

Thus switching and coordinated relabelling do not change the comparison.
Lift normalization and the cross-identification of the two five-labelled
systems remain visible inputs.  Proposition~\ref{thm:orientation-source}
makes no descent claim for them.

There is no nonzero rotation-equivariant linear comparison between harmonic
degrees three and six: they are nonisomorphic irreducible rotation modules of
dimensions seven and thirteen.  The linear comparison used here is the unique
marked \(A_5\)-intertwiner on the chosen Clebsch four-space, and it stops at
the relative marked source class
\([C_{\mathfrak m},Z_{\mathfrak m}]_{\rm or}\).
